\lecture{1}
\title{Introduction to Machine Learning}

\begin{document}

\subsection{Setting Up}

Here's the classic machine learning setup.

\begin{quote}
    \begin{itemize}
        \item There is an \textbf{input space} $\mathcal{X}$ and
        an \textbf{output space} $\mathcal{Y}$.
        \begin{itemize}
            \item The input space is often $\mathbb{R}^D$,
            where $D$ is the \textbf{feature dimension}.
            \item The output space can be, say, $\{0, 1\}$ for a classification problem,
            or $\mathbb{R}$ for a regression problem.
        \end{itemize}
        \item There is a \textbf{training set} $\{(x^{(n)}, y^{(n)})\}_{n = 1}^N$,
        where $x^{(n)} \in \mathcal{X}$ is the $n^{\text{th}}$ input, and $y^{(n)}$ is its output.
        \item There is a \textbf{decision rule} or \textbf{hypothesis}
        $h: \mathcal{X} \to \mathcal{Y}$ inferred from the training set.
        \item There is a \textbf{loss function} $L: \mathcal{Y} \times \mathcal{Y} \to \mathbb{R}$,
        where $L(a, g)$ is the penalty for guessing $g$ when the answer is $a$.
    \end{itemize}
\end{quote}

\begin{remark}
    The loss function may be asymmetric; consider, say, email spam detection.
\end{remark}

\subsection{The Optimal Hypothesis (Given the Distribution!)}

\textbf{Definition (Risk).} A decision rule $h$ has \textbf{risk} defined by
$\mathrm{risk}(h) := \mathbb{E}[L(Y, h(X))]$.

Assume that all data $(x^{(n)}, y^{(n)})$ (training or not) is drawn i.i.d. from the same distribution.
Unsurprisingly, if we know this distribution $p(x, y)$ exactly, finding the best $h$ is not hard.

\begin{theorem}{Optimal Hypothesis}
    Suppose $\mathcal{X}$ is discrete and $\mathcal{Y} = \{1, \dots, K\}$.
    Then the following minimizes risk:
    \[ h(x) = \argmin_k \sum_{j = 1}^K 
    \underbrace{L(j, k)}_{\text{cost}} \cdot 
    \underbrace{p(y = j \mid x)}_{\text{probability}}. \]
\end{theorem}

\begin{proof}
    Immediate from the definition of risk. 
    You can generalize this to the continuous case if you'd like, too. ~ $\blacksquare$
\end{proof}

The optimal $h$ takes a pretty nice form when we consider regression problems
with a particular choice of loss.
(Refer to \href{/18.650/lectures/24}{Lecture 24} from 18.650
for additional context on why the loss function is what it is.)

\begin{theorem}{Optimal Regression Hypothesis}
    Assume $\mathcal{X} = \mathbb{R}^D$, $\mathcal{Y} = \mathbb{R}$, and $L(a, g) = (a - g)^2$.
    Then an optimal decision rule is $h(x) = \mathbb{E}[Y \mid X = x]$.
\end{theorem}

\begin{proof}
    For a particular $x$, we seek to find the $k \in \mathbb{R}$
    that minimizes the following ``expected cost'':
    \[ C(k) := \int_{y \in \mathbb{R}} (y - k)^2 \cdot p(y \mid x) \, \mathrm{d}y. \]
    Any minimizing $k$ must satisfy $C'(k) = 0$, or:
    \[
    C'(k) = \int_{y \in \mathbb{R}} -2(y - k) \cdot p(y \mid x) \, \mathrm{d}y = 0
    ~ \implies ~
    \underbrace{\left(\int_{y \in \mathbb{R}} y \cdot p(y \mid x) \, \mathrm{d}y\right)}_{\mathbb{E}[Y \mid X = x]} =
    \underbrace{\left(\int_{y \in \mathbb{R}} k \cdot p(y \mid x) \, \mathrm{d}y\right)}_{k}
    \]
    Thus $k = \mathbb{E}[Y \mid X = x]$ is a critical point.
    Furthermore, $C''(k) > 0$ for all $k \in \mathbb{R}$, as shown below.
    \[ C''(k) = \int_{y \in \mathbb{R}} 2 \cdot p(y \mid x) \, \mathrm{d}y = 2 \]
    Thus $k = \mathbb{E}[Y \mid X = x]$ must be the unique minimizing value. ~ $\blacksquare$
\end{proof}

Of course, the challenge of ML is that we don't know what $p$ is\dots

\end{document}